\chapter{Experimental Setup}
\label{ch:setup}

\section{Implementation}

All code is in one Jupyter notebook, \texttt{kan\_rec.ipynb}, backed by a small
Python package, \texttt{kanrec}:

\begin{table}[H]
\centering
\caption{Project structure.}
\label{tab:files}
\small
\resizebox{\textwidth}{!}{%
\begin{tabular}{ll}
\toprule
\textbf{File} & \textbf{Role} \\
\midrule
\texttt{kan\_rec.ipynb} & Runs everything: download, preprocessing, training, tables, figures \\
\texttt{kanrec/kan.py} & KAN layer (four bases), residual block, parameter-matched MLP \\
\texttt{kanrec/data\_session.py} & YooChoose loading and GRU4Rec-protocol splits \\
\texttt{kanrec/data\_basket.py} & Instacart and Dunnhumby pipelines, product clustering \\
\texttt{kanrec/session.py} & KAN variants of GRU4Rec on the official training loop \\
\texttt{kanrec/session\_baselines.py} & Pop, S-Pop, item-kNN \\
\texttt{kanrec/basket.py} & Basket-GRU variants, frequency baselines, TIFU-KNN \\
\texttt{kanrec/metrics.py} & Session and next-basket metrics \\
\texttt{kanrec/experiments.py} & Resumable runs; one result file per model and seed \\
\texttt{kanrec/report.py} & LaTeX tables and significance tests for this report \\
\texttt{tests/} & Unit tests for the KAN layer, metrics and model wrappers \\
\texttt{third\_party/GRU4Rec\_PyTorch\_Official} & Official GRU4Rec code \cite{hidasi2018} \\
\bottomrule
\end{tabular}}
\end{table}

The official GRU4Rec code needed one compatibility change for current versions of
pandas (converting a \texttt{Series} to an array before creating a tensor); it
does not change the algorithm. The official ``compatibility mode''
initialisation, which produces float64 weights under NumPy 2, is not used; both
the reference run and all variants use the standard PyTorch initialisation of
the official model. Software: Python 3.12 and PyTorch \cite{paszke2019}.

\section{Compute Profiles}

A single variable at the top of the notebook selects the hardware profile:

\begin{itemize}
    \item \textbf{mac}: Apple M1 Pro, PyTorch MPS backend. YooChoose 1/64 with
          the official GRU4Rec size, Dunnhumby in full, a 10\% user sample of
          Instacart, one seed. The results reported in this version of the report
          come from this profile.
    \item \textbf{upes}: the UPES HPC cluster's NVIDIA Blackwell GPU. Adds
          yoochoose1/4, the full Instacart data, larger basket models and three
          seeds.
    \item \textbf{colab}: a Google Colab GPU, with intermediate sizes; results are
          written to Google Drive so that an interrupted session resumes where it
          stopped.
\end{itemize}

Every model run stores its configuration, per-epoch validation history, test
metrics, model weights, and per-prediction metric values in its own file, so an
interrupted run never repeats completed work.

\section{Hyperparameters}

\begin{table}[H]
\centering
\caption{Hyperparameters (mac profile).}
\label{tab:hyper}
\small
\resizebox{\textwidth}{!}{%
\begin{tabular}{lll}
\toprule
& \textbf{Session (GRU4Rec)} & \textbf{Next-basket (Basket-GRU)} \\
\midrule
Hidden size & 480 (official) & 512 \\
Loss & Softmax cross-entropy, $\log Q$ correction & Binary cross-entropy over all items \\
Negatives & 2048 shared, popularity$^{0.2}$ & all items \\
Optimiser & Adagrad, lr 0.07 (added parameters $10^{-3}$) & Adam, lr $3 \times 10^{-3}$, grad.\ clip 5 \\
Batch & 48 parallel sessions & 64 users \\
Dropout & 0.2 on hidden state & 0.3 on input and hidden \\
Epochs & 10, best epoch on validation & up to 30, early stopping (patience 3) \\
Model selection & Validation Recall@20 & Validation Recall@10 \\
KAN basis, grid & RBF, $G=8$ (ablation: all bases, $G \in \{4,8,12\}$) & RBF, $G=8$ \\
Seed & 42 & 42 \\
\bottomrule
\end{tabular}}
\end{table}

All variants of a model family share every hyperparameter; only the network
changes. The GRU4Rec settings are the official tuned values for YooChoose. For
the basket-GRU, the learning rate ($10^{-3}$ or $3\times10^{-3}$), batch size (16
or 64) and hidden size (512 or 1024) were chosen on the Dunnhumby
\emph{validation} users with the plain basket-GRU, before any KAN variant was
trained. The best settings (lr $3\times10^{-3}$, batch 64, 512 units) reached a
validation Recall@10 of 0.183; 1024 units were not better (0.178) and five times
slower.

\section{Training the Added Parameters}
\label{sec:stability}

The official GRU4Rec optimiser, Adagrad with learning rate 0.07, takes a first
step of roughly the full learning rate for every parameter. This is harmless
inside the GRU, where sigmoid and tanh bound every activation, but a residual
block adds an unbounded term to the hidden state that is used directly for
scoring. In a 2,000-step probe at the official model size, both the KAN head and
the MLP head diverged (training loss around 54, against 8.1 for GRU4Rec), with or
without zero initialisation. Two changes, applied identically to every KAN and
MLP variant, fix this:

\begin{itemize}
    \item the output layer of every residual branch is initialised to zero, so
          each variant starts as exactly the baseline model;
    \item parameters that are not part of the official network (the residual
          blocks and the KAN parts of the KAN cells) use Adagrad with learning
          rate $10^{-3}$, while all official parameters keep 0.07.
\end{itemize}

With these settings the KAN and MLP heads train as stably as GRU4Rec
(2,000-step training loss 8.08 and 8.09, against 8.14). Learning rates of
$3 \times 10^{-3}$ and $3 \times 10^{-4}$ gave similar losses, and $10^{-2}$ was
unstable early in training. The probe log is included in the repository.
This step matters for the interpretation of the first draft as well: an
added block that is trained with an unsuitable optimiser will lose to its
baseline whether or not it is a KAN.

\section{Evaluation Metrics}

\subsection{Session-Based}

For each prediction, every item is ranked by score. The rank of the target is
computed conservatively, counting all items whose score is at least the target's
(ties count against the model), as in the official GRU4Rec evaluation. Over the
set $T$ of all test predictions:
\begin{align}
\mathrm{Recall@}K &= \frac{1}{|T|} \sum_{t \in T} \mathbb{1}[\mathrm{rank}_t \le K], \\
\mathrm{MRR@}K &= \frac{1}{|T|} \sum_{t \in T} \frac{\mathbb{1}[\mathrm{rank}_t \le K]}{\mathrm{rank}_t}, \\
\mathrm{NDCG@}K &= \frac{1}{|T|} \sum_{t \in T} \frac{\mathbb{1}[\mathrm{rank}_t \le K]}{\log_2(\mathrm{rank}_t + 1)}.
\end{align}
With a single target per prediction, Recall@$K$ equals the hit rate HR@$K$. MRR
averages over all predictions, counting a miss as zero.

\subsection{Next-Basket}

For a user $u$ with target basket $B_u$ and ranked list $L_u$:
\begin{align}
\mathrm{Recall@}K &= \frac{|L_u^{K} \cap B_u|}{|B_u|}, \qquad
\mathrm{PHR@}K = \mathbb{1}\big[|L_u^{K} \cap B_u| > 0\big], \\
\mathrm{NDCG@}K &= \frac{\sum_{r=1}^{K} \mathbb{1}[L_{u,r} \in B_u] / \log_2(r+1)}{\sum_{r=1}^{\min(K, |B_u|)} 1/\log_2(r+1)},
\end{align}
averaged over users. Following van Maasakkers et al.\ \cite{vanmaasakkers2023},
we also report precision at the true basket size, $\mathrm{P@}|B| = |L_u^{|B_u|} \cap B_u| / |B_u|$
(equal to recall at that cut-off), recall at $2|B|$, and the average rank of the
purchased items (ties share the mid-rank). Recall@10 is further split into
\emph{repeat recall}, over target items the user had bought before, and
\emph{explore recall}, over new items \cite{li2023}.

\section{Statistical Testing}

Differences between a variant and its baseline are tested with the paired
Wilcoxon signed-rank test on per-prediction (session) or per-user (basket)
metric values, with Holm correction across all comparisons in a table. With the
large test sets used here, even small differences can be significant, so effect
sizes are reported alongside $p$-values. On the GPU profiles the mean and
standard deviation over three seeds are also reported.

\section{Decision Rule}

A KAN placement is considered helpful only if, on the \emph{validation} set, it
beats both the baseline and its parameter-matched MLP control. This rule was
fixed before the results were seen, so that the choice of what to report is not
influenced by the test set.
